\section{Introduction}
\label{sec:intro}

EEG foundation models tokenize a recording before they model it, and the
choice of what a token carries decides what the model can learn cheaply and
what it can learn only by accident. Across the current generation the
choice varies in how much of the signal survives tokenization --- magnitude
spectrograms \citep{yang2023biot,pradeepkumar2026tfm}, raw patches
\citep{wang2025cbramod,wang2024eegpt}, per-bin spectral amplitude and phase
\citep{jiang2024labram} --- but not in one respect: every token describes
one frequency band, or one patch of the broadband waveform, on its own. No
token type represents the \emph{relation} between two bands.

That relation is not an exotic quantity. Phase--amplitude coupling, the
locking of a fast band's envelope to the phase of a slow rhythm, has a
standard estimator \citep{tort2010measuring}, a physiology documented for
two decades \citep{canolty2006high,canolty2010functional}, and direct
clinical relevance to precisely the tasks foundation models are evaluated
on: it is stronger in the seizure onset zone than in unaffected tissue,
interictally as well as ictally \citep{amiri2016phase}, and periodic
epileptiform discharges are, by definition, fast activity riding a slow
rhythm. A representation from which such a quantity could in principle be
decoded is not yet one that encodes it. Magnitude tokenizers make it hard
to reach, because re-deriving phase from overlapping magnitude frames is a
phase-retrieval problem \citep{griffin1984signal} that nothing in the
pipeline performs; waveform tokenizers keep it reachable but give the
encoder no incentive to extract it, and probing studies find that the
resulting embeddings under-represent exactly the oscillatory structure
such relations live in \citep{kommineni2026aperiodic}.

This paper puts the relation into the token and asks how much a clinical
EEG model gains from having it there explicitly, and how much that depends
on the encoder it is paired with. We build a tokenizer that decomposes
each electrode into learned frequency bands, measures a coupling
statistic between them at patch resolution, and hands it to the encoder
as an interaction token alongside the band's waveform token, with
properties that fix what the token can and cannot carry. We then hold the
tokenizer fixed and vary the encoder: a small encoder built around the
token, in which bands meet only through spatial attention and no
frequency-specific structure exists; and three published foundation-model
encoders, one of which keeps a spectral branch of its own. The benefit of
the explicit coupling content is large in the first, large in the two
published encoders without a spectral pathway, and smaller and
corpus-dependent in the one with it. We read this as evidence that the
value of cross-frequency content as a token depends on the encoder's
inductive bias, and we report it as such rather than as a universal
property of the token.

\paragraph{Contributions.}
\begin{itemize}[leftmargin=*,itemsep=2pt,topsep=2pt]
\item \textbf{Cross-frequency coupling as token content.} Each band
  contributes a waveform token and an interaction token built from the
  band's measured coupling to slower bands. The interaction token is
  invariant to the phase reference of the analysis for every band that has
  a slower band to align to, preserves the band's amplitude features
  exactly, and is gated so that the coupling content can be switched off
  within one parameterization (Section~\ref{sec:tokenizer}).
\item \textbf{A small encoder built around it, with the attribution that
  goes with it.} CroFreMo folds the band rows into spatial attention and
  has no frequency sub-layer; at 2.7M parameters it is competitive with
  25--69M foundation models on twelve clinical corpora and ahead of the
  strongest reproduced baseline by $0.15$ and $0.09$ AUC-PR on the two
  seizure-detection corpora with adequate positives. With the encoder
  fixed, switching the coupling content off costs $0.136$ on CHB-MIT,
  $0.107$ on TUEV, $0.063$ on TUSZ and $0.027$ on IIIC
  (Section~\ref{sec:attribution}).
\item \textbf{Encoder-dependent benefits, measured across four
  encoders.} In place of the native tokenizer, the token lifts LaBraM and
  REVE trained from scratch by $0.15$ and $0.24$ $\kappa$ on TUEV; added to
  CBraMod's encoder next to its own spectral branch, it is attributable on
  event-classification corpora and on CHB-MIT and negative on TUSZ
  (Section~\ref{sec:content-structure}). The scope of a tokenizer claim is
  set by the encoders it is paired with, and we report it that way.
\end{itemize}

Two boundaries are stated up front. The axis-free model wins or ties on
the six corpora whose label is a paroxysmal event, a pattern, or a
recording-level abnormality, and loses to a pretrained or a larger
from-scratch model on sleep staging, cohort-level diagnosis, artifact
labelling and the smallest seizure corpus. And the cells of the axis-free
model and of the transplants are single-seed at the time of writing,
against three-seed baselines; the attribution effects are several times
the seed variance we measure with the same recipe on the same corpora,
the main-table margins on the small corpora are not, and the text says
which is which.
